% ============================================================
%  chapters/ch01-world-before-objects.tex
% ============================================================

\chapter{The World Before Objects}

\chapterepigraph{%
  The map is not the territory, but the territory is not\\
  an object either — it is a field of accessible trajectories.
}{Flyxion, \textit{Semantic Infrastructure}}

\sidenote{RSVP: see Ch.\,18}

This monograph begins not with things but with \emph{access}. The received view
in both physics and computation treats objects as primitive: particles, states,
data structures. We argue instead that \emph{accessibility landscapes} are the
more fundamental substrate, and that objects emerge as \emph{stable projections}
of constraint geometry.

The key shift is from asking \emph{what exists} to asking \emph{what is
reachable from here}.

% ── 1.1 ──────────────────────────────────────────────────────
\section{Why Trajectories Are More Fundamental Than Things}

\begin{conceptaside}{The Accessibility Inversion}
  Classical ontology asks: \emph{what things are there?} Accessibility
  ontology asks: \emph{what transitions are available to a system in state
  $q$ at time $t$?}

  This is not a merely philosophical point. In the RSVP framework, the
  scalar field $\Phifield$ encodes \emph{potential accessibility}, the
  vector field $\vfield$ encodes \emph{directed flow toward accessible
  states}, and the entropy $\Sfield$ measures the \emph{remaining freedom}
  of the system. Objects — particles, thoughts, texts — are basins in this
  landscape.
\end{conceptaside}

A trajectory through configuration space carries strictly more information than
any instantaneous state description. The gesture literature \autocite{mcneill1992}
makes this vivid: a pointing gesture cannot be decomposed into a sequence of
hand-positions without losing its meaning. The \emph{arc} is the meaning.

\begin{definition}{Trajectory}{trajectory}
  Let $\mathcal{Q}$ be a configuration space with metric $g$. A
  \emph{trajectory} is a smooth map
  \[
    \gamma : [0, T] \longrightarrow \mathcal{Q}
  \]
  satisfying the admissibility condition $\gamma(t) \in \admissible{\constraint(t)}$
  for all $t \in [0,T]$, where $\constraint(t)$ is the active constraint set
  at time $t$.
\end{definition}

\begin{remark}
  The admissibility condition is the crucial clause. It is not enough that
  $\gamma$ be smooth — it must remain \emph{within the accessible region} at
  each instant. This is what distinguishes a gesture from arbitrary hand
  motion. See Chapter~14 for the formal definition of admissibility.
\end{remark}

% ── 1.2 ──────────────────────────────────────────────────────
\section{The Map–Territory Distinction}

Alfred Korzybski's maxim is often quoted but rarely made precise. We offer
a precise formulation in terms of projection:

\begin{theorem}{Map–Territory Projection Theorem}{map-territory}
  Let $\mathcal{T}$ be a territory (a topological space with accessibility
  structure $\Acc$) and let $\mathcal{M}$ be a map (a representation space).
  Any faithful map-making procedure is a continuous projection
  \[
    \proj : \mathcal{T} \longrightarrow \mathcal{M}
  \]
  with non-trivial fibers $\fiber{\proj}{m} \neq \{*\}$ for generic $m \in \mathcal{M}$.
  Information not recovered from $\proj$ alone lives in the fiber.
\end{theorem}

\begin{proof}
  Suppose $\proj$ were injective. Then the map would encode the territory
  completely, contradicting the finite description length of any symbol system
  relative to the continuous structure of $\mathcal{T}$. By the Löwenheim–Skolem
  theorem applied to the first-order theory of $\mathcal{T}$, there exist
  non-isomorphic models sharing the same map; these constitute the fiber.
\end{proof}

\begin{example}
  GPS coordinates map Earth's surface to $\mathbb{R}^2$. The fiber over any
  coordinate pair includes altitude, subsurface structure, biological activity —
  all accessible from that location but invisible in the map.
\end{example}

% ── 1.3 ──────────────────────────────────────────────────────
\section{Constraint Before Content}

\begin{historynote}
  The priority of constraint over content has deep roots. Ibn Sina's
  distinction between \emph{mahiyyah} (quiddity) and \emph{wuj\=ud}
  (existence) anticipates it: the essence of a thing is its constraint
  structure, not its realized content. In the Latin tradition, Aquinas
  transforms this into the act/potency distinction.
\end{historynote}

In modern terms: a program's type signature constrains its behavior before any
value is computed. A physical field equation constrains accessible trajectories
before any solution is specified. Grammar constrains utterances before any word
is chosen. Content \emph{projects} from constraint.

\sidedef{$\constraint$: constraint space}

\begin{definition}{Constraint Space}{constraint-space}
  A \emph{constraint space} $(\constraint, \leq)$ is a partially ordered set
  of constraints, where $c_1 \leq c_2$ means $c_1$ is \emph{more restrictive}
  than $c_2$. The \emph{accessibility set} at $c \in \constraint$ is
  \[
    \Acc(c) \;=\; \bigl\{ q \in \mathcal{Q} \;\big|\; q \text{ satisfies } c \bigr\}
  \]
  Constraint addition is set intersection: $\Acc(c_1 \wedge c_2) = \Acc(c_1) \cap \Acc(c_2)$.
\end{definition}

% ── 3D DIAGRAM: Constraint nesting ────────────────────────────
\begin{figure}[h]
  \centering
  \begin{tikzpicture}[scale=1.0]
    % constraint-nesting-3d.tex — nested accessibility regions (no pgfmath loops)
  % Outer region Acc(c1)
  \fill[RSVPlight, opacity=0.35] (0,0) ellipse (3.0cm and 1.55cm);
  \draw[RSVPmid, thick] (0,0) ellipse (3.0cm and 1.55cm);
  \node[RSVPmid, font=\small\itshape] at (2.6, 1.2) {$\Acc(c_1)$};

  % Middle region Acc(c1 ∧ c2)
  \fill[RSVPmid!20, opacity=0.5] (0,0) ellipse (1.85cm and 0.95cm);
  \draw[RSVPdeep, thick] (0,0) ellipse (1.85cm and 0.95cm);
  \node[RSVPdeep, font=\small\itshape] at (-2.5, 0.65)
    {$\Acc(c_1 \wedge c_2)$};

  % Inner core Acc(c1 ∧ c2 ∧ c3)
  \fill[RSVPaccent, opacity=0.88] (0,0) circle (0.62cm);
  \node[white, font=\bfseries\small] at (0,0) {$\min\Phi$};
  \node[RSVPaccent, font=\scriptsize\itshape] at (0,-0.88)
    {$\Acc(c_1 \wedge c_2 \wedge c_3)$};

  % Hardcoded gradient arrows (6 directions, no \pgfmathsetmacro)
  \draw[-{Stealth[length=4pt]},RSVPaccent!80,thick,opacity=0.75] (2.7,0) -- (0.72,0);
  \draw[-{Stealth[length=4pt]},RSVPaccent!80,thick,opacity=0.75] (-2.7,0) -- (-0.72,0);
  \draw[-{Stealth[length=4pt]},RSVPaccent!80,thick,opacity=0.75] (1.35,1.1) -- (0.36,0.3);
  \draw[-{Stealth[length=4pt]},RSVPaccent!80,thick,opacity=0.75] (-1.35,1.1) -- (-0.36,0.3);
  \draw[-{Stealth[length=4pt]},RSVPaccent!80,thick,opacity=0.75] (1.35,-1.1) -- (0.36,-0.3);
  \draw[-{Stealth[length=4pt]},RSVPaccent!80,thick,opacity=0.75] (-1.35,-1.1) -- (-0.36,-0.3);

  \node[RSVPaccent!75, font=\footnotesize\itshape] at (0,-1.9)
    {$-\nabla\kappa$ (constraint descent)};

  \end{tikzpicture}
  \caption{Nested constraint spaces as concentric accessibility regions.
    Outer region: $\Acc(c_1)$. Middle: $\Acc(c_1 \wedge c_2)$. Inner core:
    $\Acc(c_1 \wedge c_2 \wedge c_3)$, the most constrained accessible set.
    The RSVP scalar field $\Phifield$ attains its minimum at the core.}
  \label{fig:constraint-nesting}
\end{figure}

% ── 1.4 ──────────────────────────────────────────────────────
\section{Accessibility Before State}

Classical state-based systems are deterministic: a state $q$ at time $t$
determines the state at $t + \delta t$ via a transition function. But this
assumes the transition function \emph{already exists} — that the dynamics are
given. In the RSVP framework we ask: where does the transition function come
from?

\begin{warning}
  State-based models smuggle accessibility assumptions into the transition
  function without acknowledging them. When the function is undefined — at
  phase transitions, at cognitive insight, at cosmological singularities —
  the state-based model collapses to silence. Accessibility geometry provides
  the right substrate for modeling these boundary cases.
\end{warning}

The RSVP answer is: transition functions \emph{are} accessibility gradients.
The coupled field equations

\begin{align}
  \partial_t \Phifield &= -\gradient \cdot \vfield + \alpha \laplacian \Phifield
    \label{eq:phi-dynamics}\\
  \partial_t \vfield &= -(\vfield \cdot \gradient)\vfield - \gradient \Phifield
    + \beta \laplacian \vfield
    \label{eq:v-dynamics}\\
  \partial_t \Sfield &= \sigma(\Phifield, \vfield) - \gamma \Sfield
    \label{eq:s-dynamics}
\end{align}

define the \emph{accessibility relaxation dynamics}: the system flows toward
states of higher admissibility under the current constraint structure.

% ── Exercises ────────────────────────────────────────────────
\section*{Exercises}

\begin{exercise}
  Let $\mathcal{Q} = \mathbb{R}^3$ and let $\constraint$ consist of a single
  linear constraint $c : \mathbf{q} \cdot \mathbf{n} = 0$ for some unit normal
  $\mathbf{n}$. Describe $\Acc(c)$ geometrically. What is $\Acc(c \wedge c')$
  when $c'$ adds a second independent constraint?
\end{exercise}

\begin{exercise}
  Formalize the GPS example from \cref{thm:map-territory} as a projection
  $\proj : \mathbb{R}^3 \to \mathbb{R}^2$. Describe the fiber over a point
  in $\mathbb{R}^2$. What additional structure (beyond the projection) would
  be needed to recover altitude information?
\end{exercise}

\begin{exercise}[title={Harder}]
  Show that the poset $(\constraint, \leq)$ with the operation
  $c_1 \wedge c_2$ forms a semilattice. Under what conditions does it form
  a complete lattice? What is the interpretation of the top element
  $\top \in \constraint$ in terms of $\Acc$?
\end{exercise}
